\section{Experiments}
\label{sec:experiments}

We evaluate generational depth estimation on four model configurations:
Pythia-1.4B and OLMo-1B in continuation mode, and Qwen2.5-7B-Instruct in both rewrite and continuation modes.
All experiments use the OBD-LightGBM classifier with the 12-dimensional feature set described in Section~3.
Each configuration uses a balanced corpus of 500 texts per depth level (generations 0--5).
We organize results around four questions:
(1)~Does $\Kempirical = 2$ hold across configurations?
(2)~Do depth classifiers transfer across models?
(3)~Is the feature degradation ordering model-independent?
(4)~How does generation mode affect detectability?

%% ====================================================================
\subsection{$\Kstar = 2$ Across All Configurations}
\label{sec:kstar_boundary}

Table~\ref{tab:kstar_ci} reports pairwise $\auc$ for adjacent depth pairs
across all four configurations, with bootstrap 95\% confidence intervals
(1{,}000 resamples) and permutation-test $p$-values for the critical 1v2
and 2v3 boundaries.

\begin{table}[t]
\centering
\small
\caption{Pairwise $\auc$ for adjacent depth pairs.  Bootstrap 95\% CIs
and permutation $p$-values shown for the $\Kstar$ boundary rows (1v2,
2v3).  $\Kempirical = 2$ requires 1v2 $\text{CI}_\text{low} > \theta$
and 2v3 $\text{CI}_\text{high} < \theta$ with $\theta = 0.60$.
All 0v1 permutation $p = 0.002$.}
\label{tab:kstar_ci}
\begin{tabular}{lcccc}
\toprule
 & Qwen-rw & Pythia & OLMo & Qwen-cont \\
\midrule
0v1 & 0.9956 & 0.9932 & 0.9970 & 0.9972 \\[4pt]
1v2 & 0.6955 & 0.6539 & 0.6145 & 0.6112 \\
{\footnotesize CI}
  & {\footnotesize [.6859,\,.7041]}
  & {\footnotesize [.6454,\,.6630]}
  & {\footnotesize [.6045,\,.6247]}
  & {\footnotesize [.6015,\,.6213]} \\
{\footnotesize $p$}
  & {\footnotesize .002}
  & {\footnotesize .002}
  & {\footnotesize .002}
  & {\footnotesize .002} \\[4pt]
2v3 & 0.5893 & 0.5514 & 0.5496 & 0.5609 \\
{\footnotesize CI}
  & {\footnotesize [.5795,\,.5991]}
  & {\footnotesize [.5417,\,.5628]}
  & {\footnotesize [.5399,\,.5596]}
  & {\footnotesize [.5511,\,.5712]} \\
{\footnotesize $p$}
  & {\footnotesize .002}
  & {\footnotesize .002}
  & {\footnotesize .002}
  & {\footnotesize .002} \\[4pt]
3v4 & 0.5435 & 0.5152 & 0.5091 & 0.5409 \\
4v5 & 0.5340 & 0.5136 & 0.5045 & 0.5141 \\
\bottomrule
\end{tabular}
\end{table}

The pattern is uniform across all four configurations:
0v1 discrimination is near-perfect ($\auc > 0.99$),
followed by a sharp drop at 1v2 ($\auc$ range: 0.6112--0.6955),
with 2v3 falling below the decision threshold $\theta = 0.60$
($\auc$ range: 0.5496--0.5893).
All subsequent pairs (3v4, 4v5) converge toward chance.
Figure~\ref{fig:auc_decay} visualizes this decay.

\begin{figure}[t]
\centering
\includegraphics[width=\columnwidth]{figures/fig2_auc_decay.pdf}
\caption{Pairwise $\auc$ as a function of depth pair across four model
configurations.  All configurations show a sharp drop between 0v1 and
1v2, with 2v3 and deeper pairs falling below the decision threshold
$\theta = 0.60$ (dashed line).  Error bars: bootstrap 95\% CIs.}
\label{fig:auc_decay}
\end{figure}

For the $\Kstar$ determination, all four 1v2 CI lower bounds exceed
$\theta$: Qwen-rw ($\text{CI}_\text{low} = 0.6859$), Pythia (0.6454),
OLMo (0.6045), and Qwen-cont (0.6015).
Conversely, all four 2v3 CI upper bounds fall below $\theta$:
Qwen-rw ($\text{CI}_\text{high} = 0.5991$), Qwen-cont (0.5712),
Pythia (0.5628), and OLMo (0.5596).
Permutation tests confirm $p = 0.002$ for all eight boundary comparisons.
Together these establish $\Kempirical = 2$ across all tested configurations.

The most marginal case is Qwen-cont, where the 1v2 CI lower bound
(0.6015) exceeds $\theta$ by only 0.0015.  This narrow margin
indicates that $\Kstar = 2$ is robust but not trivially so---continuation-mode
generation produces weaker depth artifacts than rewrite mode, pushing the
boundary closer to the detection threshold.

The absolute 1v2 $\auc$ varies across models:
Qwen-rw (0.6955) $>$ Pythia (0.6539) $>$ OLMo (0.6145) $\approx$ Qwen-cont (0.6112).
Notably, the spread due to generation mode within a single model
(Qwen-rw vs.\ Qwen-cont: $\Delta = 8.4$\,pp) exceeds the spread
across models in the same mode
(Pythia vs.\ OLMo: $\Delta = 3.9$\,pp),
suggesting that generation protocol is a stronger determinant
of detectability than model identity.
We examine both factors in Sections~\ref{sec:transfer}
and~\ref{sec:generation_mode}.

Beyond the $\Kstar$ boundary, depth discrimination deteriorates
rapidly.
OLMo shows nonsignificant permutation $p$-values at 3v4
($p = 0.064$) and 4v5 ($p = 0.232$), indicating that depth
classification beyond generation~2 is indistinguishable from
random for this model.
Pythia retains marginal effects at these depths
($p = 0.006$ and $p = 0.008$, respectively),
while Qwen-cont shows weak but detectable discrimination through
3v4 ($p = 0.002$) that fades at 4v5 ($p = 0.006$).

%% ====================================================================
\subsection{Cross-Model Transfer}
\label{sec:transfer}

To assess whether depth signals are properties of the text itself
rather than artifacts of a specific generator, we train classifiers
on one model's outputs and evaluate on another's.
Table~\ref{tab:transfer_matrix} shows the full
$4 \times 4$ transfer accuracy matrix for 6-class ordinal depth
classification.

\begin{table}[t]
\centering
\small
\caption{Cross-model transfer accuracy (6-class ordinal
classification).  Diagonal entries (bold): within-model accuracy.
Off-diagonal: cross-model transfer.  Retention is computed as
off-diagonal accuracy divided by the diagonal entry of the training
model.}
\label{tab:transfer_matrix}
\begin{tabular}{lcccc}
\toprule
Train $\downarrow$ / Test $\rightarrow$
  & Qwen-rw & Pythia & OLMo & Qwen-cont \\
\midrule
Qwen-rw
  & \textbf{0.4229} & 0.2590 & 0.2653 & 0.3466 \\
Pythia
  & 0.2718 & \textbf{0.3856} & 0.3531 & 0.3287 \\
OLMo
  & 0.2768 & 0.3450 & \textbf{0.3822} & 0.3244 \\
Qwen-cont
  & 0.3440 & 0.3284 & 0.3112 & \textbf{0.3875} \\
\bottomrule
\end{tabular}
\end{table}

\begin{figure}[t]
\centering
\includegraphics[width=0.85\columnwidth]{figures/fig3_transfer_heatmap.pdf}
\caption{Cross-model transfer heatmap (6-class ordinal accuracy).
Rows: training model; columns: test model.
Pythia$\leftrightarrow$OLMo (both continuation mode) shows high
mutual transfer (91--92\% accuracy retention), while Qwen-rewrite
transfers poorly to continuation-mode models (61--63\%).
Diagonal entries are within-model baselines.}
\label{fig:transfer_heatmap}
\end{figure}

Two principal findings emerge.
First, transfer between continuation-mode models is high.
Pythia$\to$OLMo retains 91--92\% of within-model accuracy, and
the reverse direction is comparable
(85--86\% retention by Kendall $\tau$).
The accuracy-based retention of 91--92\% is complemented by the
$\tau$-based retention of 85--86\%, indicating that not only
classification accuracy but also the relative depth ordering of
predictions is largely preserved.
Qwen-cont$\leftrightarrow$Pythia also transfers well
(83--85\% accuracy retention),
and Qwen-cont$\leftrightarrow$OLMo achieves 80--85\%.
These retention rates indicate that a substantial fraction of
the depth signal---on the order of 70--85\%---is text-intrinsic
rather than tied to a specific generator.

Second, Qwen-rewrite exhibits catastrophic transfer failure.
Classifiers trained on Qwen-rewrite retain only 61--63\% of
within-model accuracy when evaluated on Pythia or OLMo
(45--50\% by $\tau$).
The reverse direction performs comparably poorly.
This confirms that rewrite-mode classifiers rely on generation
artifacts that do not generalize across models, in contrast
to Qwen-cont, which achieves 80--85\% retention to the same targets.

At the 0v1 boundary, cross-model pairwise $\auc$ exceeds 0.96
for all model pairs, confirming that the human-vs.-synthetic
distinction is universally transferable regardless of generator
or mode.
The transfer degradation concentrates at deeper boundaries
(1v2 and beyond), where model- and mode-specific factors dominate.

The aggregate picture is one of modest cross-model loss:
mean pairwise $\auc$ is 0.6099 within-model versus 0.5905
cross-model.
This gap is driven primarily by Qwen-rewrite; among
continuation-mode pairs the within/cross difference narrows
further.

%% ====================================================================
\subsection{Feature Degradation Ordering}
\label{sec:feature_degradation}

If depth signals are text-intrinsic, different generators should
rely on the same features in the same order.
We test this by computing single-feature $\auc$ rankings
independently for each model and comparing them via rank
correlation.

\begin{table}[t]
\centering
\small
\caption{Top-3 features by single-feature $\auc$ ranking for each
model configuration.
Features shared across all three continuation-mode models are in bold.
Qwen-rewrite's top feature (\texttt{freq\_entropy}) is mode-specific.
Spearman $\rho$ (continuation-mode pairs): Pythia--OLMo 0.95,
Pythia--Qwen-cont 0.9143, OLMo--Qwen-cont 0.85.}
\label{tab:top_features}
\begin{tabular}{cllll}
\toprule
Rank & Pythia & OLMo & Qwen-rw & Qwen-cont \\
\midrule
1
  & \textbf{\texttt{surp\_d1\_std}}
  & \textbf{\texttt{surp\_std}}
  & \texttt{freq\_entropy}
  & \textbf{\texttt{surp\_d1\_std}} \\
2
  & \textbf{\texttt{surp\_d2\_std}}
  & \textbf{\texttt{surp\_d1\_std}}
  & \textbf{\texttt{surp\_std}}
  & \textbf{\texttt{surp\_std}} \\
3
  & \textbf{\texttt{surp\_std}}
  & \textbf{\texttt{surp\_d2\_std}}
  & \textbf{\texttt{surp\_d1\_std}}
  & \textbf{\texttt{surp\_d2\_std}} \\
\bottomrule
\end{tabular}
\end{table}

Among continuation-mode models, single-feature $\auc$ rankings
are highly correlated:
Pythia vs.\ OLMo $\rho = 0.95$,
Pythia vs.\ Qwen-cont $\rho = 0.9143$,
OLMo vs.\ Qwen-cont $\rho = 0.85$
(all $\rho \geq 0.85$).
The top-3 features---\texttt{surp\_d1\_std}, \texttt{surp\_std},
and \texttt{surp\_d2\_std}---are shared (up to ordering) across
all three continuation-mode configurations
(Table~\ref{tab:top_features}).
All three measure variability of surprisal or its derivatives,
capturing the \emph{heterogeneity} of token-level predictability
across the text.
That these features degrade in a shared order across generators
suggests a common mechanism: each generation step smooths the
token-level surprisal landscape, reducing its variance in a
model-independent fashion.

Qwen-rewrite diverges: its top-ranked feature is
\texttt{freq\_entropy}, a lexical diversity measure absent from
the continuation-mode top-3 for any model.
This corroborates the transfer results in
Section~\ref{sec:transfer}---rewrite-mode classifiers learn
partly from instruction-following artifacts rather than the
shared surprisal-based signals.

Kendall $\tau$ provides a complementary, more conservative
perspective:
Pythia vs.\ OLMo $\tau = 0.733$,
Pythia vs.\ Qwen-cont $\tau = 0.695$,
OLMo vs.\ Qwen-cont $\tau = 0.543$.
Cross-mode correlations are lower:
Pythia vs.\ Qwen-rw $\tau = 0.467$,
OLMo vs.\ Qwen-rw $\tau = 0.429$.

\paragraph{Limitation.}
The high Spearman $\rho$ characterizes the \emph{aggregate}
feature ranking, but at the depth-pair level the agreement
weakens: $\tau \approx 0.37$ at the 1v2 boundary.
This means that while the overall importance ordering is shared,
the relative contribution of individual features at the
discrimination boundary---where precision matters most---varies
across models.
We view this as a genuine limitation of the text-intrinsic
interpretation and a direction for future investigation.

%% ====================================================================
\subsection{Generation Mode Effect}
\label{sec:generation_mode}

Qwen2.5-7B-Instruct provides a controlled comparison of generation
modes: rewrite (instruction-based paraphrasing) and continuation
(prompt truncation), as defined in Section~3.
Both share the same model weights, isolating the effect of
generation protocol on depth detectability.

Rewrite mode yields higher 1v2 $\auc$ than continuation:
0.6955 vs.\ 0.6112, a difference of $+$8.4 percentage points.
The advantage persists at 2v3
(0.5893 vs.\ 0.5609, $\Delta = +$2.8\,pp)
and narrows at 3v4
(0.5435 vs.\ 0.5409, $\Delta = +$0.3\,pp),
with both modes converging toward chance at deeper levels.
The amplified detectability under rewrite mode is consistent
with instruction-following introducing systematic deviations
from the source text's style---additional artifacts that are
discriminable but mode-specific rather than text-intrinsic.

Transfer between modes is asymmetric.
A classifier trained on Qwen-cont and evaluated on Qwen-rw
retains 89\% of within-model accuracy
(0.3440 vs.\ 0.3875 diagonal),
whereas the reverse direction---Qwen-rw trained,
Qwen-cont tested---retains 82\%
(0.3466 vs.\ 0.4229 diagonal).
The direction of the asymmetry is intuitive:
continuation-mode signals form a subset of the information
available in rewrite-mode data
(which adds instruction-following artifacts on top),
so a continuation-trained classifier can partially recognize
rewrite-mode outputs but not vice versa.

This asymmetry is mirrored in the feature rankings.
Qwen-rewrite's top feature (\texttt{freq\_entropy}) is
mode-specific and ranks outside the top-3 for every
continuation-mode model.
Qwen-cont's top feature (\texttt{surp\_d1\_std}) is shared
with Pythia and OLMo, confirming that continuation mode
captures primarily text-intrinsic depth signals.

The practical implication is a detectability--generalizability
tradeoff.
Rewrite-mode classifiers achieve higher absolute $\auc$ but
transfer poorly across models
(Qwen-rw$\to$Pythia/OLMo: 61--63\% retention).
Continuation-mode classifiers sacrifice 8.4\,pp of within-model
$\auc$ but achieve robust cross-model transfer
(Qwen-cont$\to$Pythia: 83--85\%).
For deployment scenarios requiring model-agnostic depth
estimation, continuation-mode features provide a more reliable
foundation despite lower raw discriminability.
